\uuid{8RMU}
\exo7id{4820}
\titre{exo7 4820}
\auteur{quercia}
\organisation{exo7}
\datecreate{2010-03-16}
\isIndication{false}
\isCorrection{false}
\chapitre{Topologie}
\sousChapitre{Topologie des espaces vectoriels normés}
\cpge{mpsi-mp2i.nombres-reels-suites.b, pcsi.nombres-reels-suites.a, ptsi.nombres-reels-suites.a, mpsi-mp2i.nombres-reels-suites.d, pcsi.nombres-reels-suites.c, ptsi.nombres-reels-suites.c, bcpst1.suites-reelles, mp-mpi.topologie-espaces-normes.a, pc.espaces-vectoriels-normes.a, psi.espaces-vectoriels-normes.a, mp-mpi.topologie-espaces-normes.b, pc.espaces-vectoriels-normes.b, psi.espaces-vectoriels-normes.b, mp-mpi.topologie-espaces-normes.j, pc.espaces-vectoriels-normes.g, psi.espaces-vectoriels-normes.g}
\module{Analyse}
\niveau{L2}
\difficulte{}

\contenu{
\texte{
Soit $E$ un evn de dimension finie et $u \in \mathcal{L}(E)$. On pose $x_n = \|u^n\|$.
}
\begin{enumerate}
    \item \question{Montrer que $\rho = \inf\{\sqrt[n]{x_n},\ n\in \N\}$ est ind{\'e}pendant de la norme
    choisie sur $E$.}
    \item \question{En utilisant l'in{\'e}galit{\'e} : $x_{p+q} \le x_p x_q$, montrer que
    la suite $(\sqrt[n] {x_n}\,)$ converge vers $\rho$.}
\end{enumerate}
}
